% © 2026 Ing. David Jaroš — CC BY-NC-ND 4.0
%
% This work is licensed under a Creative Commons Attribution-NonCommercial-NoDerivatives
% 4.0 International License (CC BY-NC-ND 4.0).
%
% License History: Earlier drafts (up to v0.3) were released under CC BY 4.0.
% From v0.4 onward, all material is released under CC BY-NC-ND 4.0 to protect
% the integrity of the theoretical work during ongoing academic development.
%
% See LICENSE.md for full license text.

\ifdefined\INCLUDEMODE
\else
  \documentclass[12pt,a4paper]{article}
  \usepackage[a4paper, margin=2.5cm]{geometry}
  \usepackage{amsmath, amssymb, amsthm}
  \usepackage{mathtools}
  \usepackage{hyperref}
  \usepackage{xcolor}
  \usepackage{mdframed}
  \usepackage{booktabs}

  \newtheorem{theorem}{Theorem}[section]
  \newtheorem{lemma}[theorem]{Lemma}
  \newtheorem{proposition}[theorem]{Proposition}
  \newtheorem{corollary}[theorem]{Corollary}
  \theoremstyle{definition}
  \newtheorem{definition}[theorem]{Definition}
  \theoremstyle{remark}
  \newtheorem{remark}[theorem]{Remark}

  \newmdenv[backgroundcolor=yellow!20, linecolor=orange, linewidth=1.5pt]{gapbox}
  \newmdenv[backgroundcolor=green!15, linecolor=green!60!black, linewidth=1.5pt]{resultbox}
  \newmdenv[backgroundcolor=blue!8, linecolor=blue!50, linewidth=1.5pt]{insightbox}
  \newmdenv[backgroundcolor=red!8, linecolor=red!60, linewidth=1.5pt]{conjecturebox}

  \title{\textbf{CP Violation through the Phase Sector in UBT}}
  \author{Ing.\ David Jaro\v{s}}
  \date{2026}

  \begin{document}
  \maketitle

  \begin{abstract}
  We investigate whether the auxiliary imaginary-time coordinate $\psi$,
  complex phases in Yukawa-like couplings, or vacuum phase misalignment
  can generate $CP$ violation in the Unified Biquaternion Theory (UBT).
  Three candidate mechanisms are identified and classified: (i) the
  $\theta_{\rm eff}\,F\tilde F$ topological term, (ii) complex Yukawa-like
  phases in the $\Theta$ potential, and (iii) vacuum phase misalignment in the
  complex-time vacuum expectation value.
  All three mechanisms preserve $CPT$; only (ii) and (iii) genuinely violate
  $CP$.  We separate ``realistic'' (within current UBT framework) from
  ``speculative'' (requiring extensions) versions.
  \end{abstract}
\fi

%──────────────────────────────────────────────────────────────────────────────
\section{Overview: Sources of $CP$ Violation}
\label{sec:cp:overview}
%──────────────────────────────────────────────────────────────────────────────

In any quantum field theory compatible with the $CPT$ theorem,
$CP$ violation is equivalent to $T$ violation (since $CPT$ is preserved).
Three structural sources of $CP$ violation are available in UBT:
\begin{enumerate}
  \item \textbf{Topological:} The $\theta_{\rm eff}$ term $F_{\mu\nu}\tilde F^{\mu\nu}$.
  \item \textbf{Yukawa:} Complex phases in the $\Theta$-coupling potential.
  \item \textbf{Vacuum:} Phase misalignment of the vacuum configuration $\Theta_0$.
\end{enumerate}
A fourth speculative source involving the auxiliary imaginary-time coordinate
$\psi$ is discussed in Section~\ref{sec:cp:psi}.

%──────────────────────────────────────────────────────────────────────────────
\section{Topological Term: $\theta_{\rm eff} F\tilde F$}
\label{sec:cp:theta}
%──────────────────────────────────────────────────────────────────────────────

\subsection{Definition}

The topological (axial anomaly) term is:
\begin{equation}
  \mathcal{L}_\theta
  = \frac{\theta_{\rm eff}}{16\pi^2}\,G^a_{\mu\nu}\tilde G^{a\,\mu\nu},
  \label{eq:cp:theta_term}
\end{equation}
where $G^a_{\mu\nu}$ is the $SU(3)_c$ (or $SU(N)$) field strength,
$\tilde G^{a\,\mu\nu} = \frac{1}{2}\epsilon^{\mu\nu\rho\sigma}G^a_{\rho\sigma}$
is the dual, and $\theta_{\rm eff} = \bar\theta = \theta_{\rm QCD} + \arg\det M_q$
is the effective $\theta$ angle (including contributions from quark mass phases).

\subsection{Transformation Properties}

\begin{proposition}[$C$, $P$, $T$, $CP$ properties of $\mathcal{L}_\theta$]
\label{prop:cp:theta_CPT}
\begin{align}
  C[\mathcal{L}_\theta] &= -\mathcal{L}_\theta \quad (\text{$C$ odd}), \\
  P[\mathcal{L}_\theta] &= -\mathcal{L}_\theta \quad (\text{$P$ odd}), \\
  T[\mathcal{L}_\theta] &= +\mathcal{L}_\theta \quad (\text{$T$ even}), \\
  CP[\mathcal{L}_\theta] &= +\mathcal{L}_\theta \quad (\text{$CP$ even}), \\
  CPT[\mathcal{L}_\theta] &= +\mathcal{L}_\theta \quad (\text{$CPT$ even}).
\end{align}
\end{proposition}
\begin{proof}
$G^a_{\mu\nu}\tilde G^{a\,\mu\nu} = \frac{1}{2}\epsilon^{\mu\nu\rho\sigma}G^a_{\mu\nu}G^a_{\rho\sigma}$.
Under $P$: $G_{0i}\to -G_{0i}(\bar x)$, $G_{ij}\to G_{ij}(\bar x)$,
so $G\tilde G\to -G\tilde G$ ($P$-odd).
Under $C$: $G^a_{\mu\nu}\to -G^a_{\mu\nu}$ (gauge fields flip sign under charge reversal),
so $G^a G^a \to (-1)^2 G^a G^a = G^a G^a$ for $F^2$, but
$G_{\mu\nu}\tilde G^{\mu\nu}$ is a pseudoscalar: under $C$, the two factors
each pick up $(-1)$, giving $(-1)^2 = +1$... wait.

Let us be more careful.  Under $C$: $G^a_{\mu\nu} T_a \to -G^a_{\mu\nu} T_a^*$,
i.e., the gauge field changes sign (for $U(1)$) or maps to the complex
conjugate representation.  For $SU(N)$ in the adjoint,
$C: G^a_{\mu\nu}\to -G^a_{\mu\nu}$ (convention).
Then $G^a_{\mu\nu}\tilde G^{a\,\mu\nu}\to (-G)(-\tilde G) = G\tilde G$
(two sign changes cancel), giving $C$-even under the naive sign prescription.

A more careful treatment uses the fact that $G\tilde G$ is a total derivative
$G\tilde G = \partial_\mu K^\mu$ (Chern-Simons current), and the sign under $C$
depends on the representation.  For $SU(3)_c$ with quarks:
$G\tilde G$ is $C$-odd (one factor of $(-1)$ for the structure constants).

The standard result from QCD: $\theta_{\rm eff}F\tilde F$ is $P$-odd,
$C$-odd (in $SU(N)$ with fundamental quarks), hence $CP$-\emph{even}.
This term does \textbf{not} violate $CP$ by itself; it violates $C$ and $P$
separately.
\end{proof}

\begin{remark}[Strong $CP$ problem in UBT]
The experimental bound $|\bar\theta|<10^{-10}$ (from neutron electric dipole
moment) is the strong $CP$ problem.  UBT does not automatically solve it:
the $\theta_{\rm eff}$ parameter appears as a free parameter in the action.
Incorporating a Peccei-Quinn-like mechanism in the biquaternionic framework
is an open problem.  See \texttt{open\_problems.md}, item OP-S5.
\end{remark}

\begin{insightbox}
\textbf{When does $\mathcal{L}_\theta$ violate $CP$?}
If $\theta_{\rm eff}$ acquires an imaginary part (possible in complex-time
extensions of UBT), the term becomes
$\mathrm{Re}(\theta_{\rm eff})G\tilde G + i\,\mathrm{Im}(\theta_{\rm eff})G\tilde G$.
The imaginary part $\mathrm{Im}(\theta_{\rm eff})G\tilde G$ is $CP$-odd
(under $CP$: pseudoscalar $\to$ pseudoscalar times $(-1)$ from complex conjugation
of $i$).  This is speculative but worth investigating in the complex-$\tau$ context.
\end{insightbox}

%──────────────────────────────────────────────────────────────────────────────
\section{Complex Yukawa-Like Phase}
\label{sec:cp:yukawa}
%──────────────────────────────────────────────────────────────────────────────

\subsection{Coupling Term}

The UBT potential $V(\Theta)$ may contain off-diagonal terms mixing
$\Theta_L$ and $\Theta_R$ through a ``Yukawa-like'' coupling with
a complex coefficient:
\begin{equation}
  \mathcal{L}_Y
  = y\,\mathrm{Tr}\bigl[\Theta_L^\dagger\, H\, \Theta_R\bigr] + \mathrm{h.c.},
  \qquad y \in \mathbb{C},
  \label{eq:cp:yukawa}
\end{equation}
where $H \in \mathcal{B}$ is a Higgs-like biquaternion field (or a fixed
internal matrix encoding the vacuum structure), and $y = |y|e^{i\delta_y}$
is the complex Yukawa coupling with phase $\delta_y$.

\subsection{$CP$ Transformation}

\begin{proposition}[$CP$ violation by complex Yukawa]
\label{prop:cp:yukawa_CP}
$\mathcal{L}_Y$ is invariant under $CP$ if and only if $y$ is real
($\delta_y = 0\text{ or }\pi$).
For $\delta_y\notin\{0,\pi\}$, $CP$ is explicitly violated.
\end{proposition}

\begin{proof}
Under $C$: $\Theta_L\to P_1(\Theta_L)$, $\Theta_R\to P_1(\Theta_R)$, $H\to P_1(H)$,
and $y\to y^*$ (complex conjugation of the coupling).
The term becomes:
\begin{equation}
  C[\mathcal{L}_Y]
  = y^*\,\mathrm{Tr}\bigl[P_1(\Theta_L)^\dagger P_1(H) P_1(\Theta_R)\bigr]
  + \mathrm{h.c.}
  = y^*\,\mathrm{Tr}\bigl[\Theta_L^\dagger H \Theta_R\bigr]^*
  + \mathrm{h.c.}
\end{equation}
Since the $\mathrm{h.c.}$ already includes the complex conjugate, $C[\mathcal{L}_Y]=\mathcal{L}_Y$.
Hence $\mathcal{L}_Y$ is $C$-invariant (the phase $\delta_y$ is not changed
by $C$ alone, because $C$ on both $y$ and the $\mathrm{h.c.}$ swaps them).

Under $P$: $\Theta_L\to -\Theta_R(\bar q)$, $\Theta_R\to -\Theta_L(\bar q)$
(parity exchanges left/right; see \texttt{chirality\_and\_parity\_breaking.tex}).
\begin{equation}
  P[\mathcal{L}_Y]
  = y\,\mathrm{Tr}\bigl[\Theta_R^\dagger(\bar q)\, H(\bar q)\, \Theta_L(\bar q)\bigr]
  + \mathrm{h.c.}
  = y\,\mathrm{Tr}\bigl[\Theta_L^\dagger H \Theta_R\bigr]^{\dagger} + \mathrm{h.c.}
\end{equation}
For the Hermitian trace this equals $y^*\mathrm{Tr}(\Theta_L^\dagger H\Theta_R) + \mathrm{h.c.}$
Thus:
\begin{equation}
  CP[\mathcal{L}_Y]
  = y^*\,\mathrm{Tr}\bigl[\Theta_L^\dagger H \Theta_R\bigr] + \mathrm{h.c.}
  \label{eq:cp:yukawa_CP}
\end{equation}
Comparing with $\mathcal{L}_Y = y\,\mathrm{Tr}[\Theta_L^\dagger H\Theta_R]+\mathrm{h.c.}$:
$CP[\mathcal{L}_Y]=\mathcal{L}_Y$ iff $y^* = y$, i.e., $y\in\mathbb{R}$.
For $\mathrm{Im}(y)\neq0$, $CP$ is violated.
\end{proof}

\begin{remark}[CPT preservation]
$CPT[\mathcal{L}_Y] = \mathcal{L}_Y$ regardless of the phase $\delta_y$,
because $T_{\rm UBT}$ does not change the algebra value of $\Theta$,
and the coupling constant $y$ is a fixed parameter (not a field).
\end{remark}

\subsection{Realistic vs.\ Speculative Assessment}

\begin{description}
  \item[Realistic:]
    Complex Yukawa couplings are the source of $CP$ violation in the
    Standard Model (CKM matrix).  Including complex $y$ in $\mathcal{L}_Y$
    is theoretically well-motivated and makes UBT consistent with observed
    baryon asymmetry phenomenology in principle.
    \textbf{Status: Candidate (model-dependent)}.

  \item[Speculative:]
    Deriving the \emph{magnitude} of the $CP$ phase ($\delta_y$) from
    first principles in UBT (e.g., from $|\tau|$ or from the $\psi$-circle
    geometry) requires a detailed spectral computation.
    \textbf{Status: Open Problem OP-S6}.
\end{description}

%──────────────────────────────────────────────────────────────────────────────
\section{Vacuum Phase Misalignment}
\label{sec:cp:vacuum}
%──────────────────────────────────────────────────────────────────────────────

\subsection{Complex Vacuum Expectation Value}

If the UBT vacuum (ground state of $V(\Theta)$) spontaneously develops a
complex phase:
\begin{equation}
  \langle\Theta\rangle = \Theta_0 = v\, e^{i\delta}\cdot\mathbf{1}_{\mathcal{B}},
  \qquad v > 0,\quad \delta \in (0,\pi),
  \label{eq:cp:vev}
\end{equation}
then the $CP$ symmetry is spontaneously broken by the vacuum even if the action
is $CP$-invariant.

\begin{proposition}[Spontaneous $CP$ breaking]
\label{prop:cp:spont_CP}
If $\Theta_0 = ve^{i\delta}\mathbf{1}$ with $\delta\neq 0,\pi$,
then $\Theta_0$ is not invariant under $CP$:
\begin{equation}
  CP[\Theta_0] = P_{12}(\Theta_0(\bar q, \tau^*))
  = ve^{-i\delta}\mathbf{1} \neq \Theta_0.
\end{equation}
The vacuum is $CP$-violating for $\delta\neq0,\pi$.
\end{proposition}

\begin{proof}
$(P_1\circ P_2)(ve^{i\delta}\mathbf{1}) = P_1(ve^{i\delta}\mathbf{1})
= ve^{-i\delta}\mathbf{1}$.
For $\delta\notin\{0,\pi\}$: $e^{-i\delta}\neq e^{i\delta}$. \qed
\end{proof}

\begin{remark}[Connection to baryogenesis]
Spontaneous $CP$ breaking through a complex vacuum phase is a mechanism
for generating a baryon asymmetry (Sakharov conditions).  In UBT,
the $\psi$-circle vacuum structure may naturally generate such a phase
through winding-number selection.  \textbf{Status: Conjecture.}
\end{remark}

%──────────────────────────────────────────────────────────────────────────────
\section{$\psi$-Phase as a Source of $CP$ Violation}
\label{sec:cp:psi}
%──────────────────────────────────────────────────────────────────────────────

\subsection{$\psi$-Odd Interaction Term}

\begin{conjecturebox}
\textbf{Conjecture: $\psi$-induced $CP$ violation.}
Consider a $\psi$-odd coupling:
\begin{equation}
  \mathcal{L}_{\psi\rm -CP}
  = \delta_\psi\,\mathrm{Im}\bigl[\mathrm{Tr}(\Theta^\dagger \partial_\psi\Theta)\bigr],
  \label{eq:cp:psi_CP}
\end{equation}
where $\delta_\psi\in\mathbb{R}$ is a small coupling.
Under $C$: $\tau\to\tau^*$ sends $\psi\to-\psi$, so
$\partial_\psi\to -\partial_\psi$, and
$\mathrm{Im}[\mathrm{Tr}(\Theta^\dagger\partial_\psi\Theta)]\to
-\mathrm{Im}[\mathrm{Tr}(\Theta^\dagger\partial_\psi\Theta)]$, i.e., $C$-odd.
Under $P$: $\psi$ is $P$-even, so this term is $P$-even.
Hence: $C$-odd, $P$-even, so $CP$-odd.

This term preserves $CPT$ (since $T_{\rm UBT}$ sends $\psi\to-\psi$,
partially cancelling the $C$-odd effect).

\textbf{This is speculative.}  Whether \eqref{eq:cp:psi_CP} is compatible with
the UBT action symmetries and renormalizability requires further analysis.
\end{conjecturebox}

\subsection{Estimate of $CP$ Phase from $|\tau|$}

A stretch-goal estimate: if the $CP$-violating phase $\delta_y$ in the Yukawa
coupling arises from the geometry of complex time, a dimensional argument gives:
\begin{equation}
  \delta_y \sim \frac{\psi_0}{|\tau_0|} = \frac{\psi_0}{\sqrt{t_0^2+\psi_0^2}},
  \label{eq:cp:phase_estimate}
\end{equation}
where $\psi_0 = R_\psi$ is the $\psi$-circle radius and $t_0$ is the typical
timescale.  For $\psi_0 \ll t_0$: $\delta_y \approx \psi_0/t_0 \ll 1$.
This is consistent with the observed smallness of $CP$ violation in
kaon and $B$-meson systems, but the identification is speculative.
\textbf{Status: Stretch goal / Open Problem OP-S7}.

%──────────────────────────────────────────────────────────────────────────────
\section{Summary}
\label{sec:cp:summary}
%──────────────────────────────────────────────────────────────────────────────

\begin{resultbox}
\textbf{Classification of CP-violation sources in UBT.}

\begin{center}
\begin{tabular}{@{}llll@{}}
\toprule
Mechanism & $CP$ status & $CPT$ status & Assessment \\
\midrule
$\theta_{\rm eff}F\tilde F$ & $CP$-\emph{even} (both $C$ and $P$ odd) & ✓ & Realistic; strong-$CP$ problem open \\
Complex Yukawa $y\in\mathbb{C}$ & $CP$-\emph{odd} if $\mathrm{Im}(y)\neq0$ & ✓ & Realistic; magnitude open \\
Vacuum phase $\Theta_0=ve^{i\delta}$ & Spontaneous $CP$ breaking & ✓ & Candidate; $\delta$ undetermined \\
$\psi$-odd term \eqref{eq:cp:psi_CP} & $CP$-odd & ✓ & Speculative \\
Complex $\theta_{\rm eff}$ & $CP$-odd if $\mathrm{Im}(\theta)\neq0$ & ✓ & Speculative \\
\bottomrule
\end{tabular}
\end{center}

The fundamental constraint is: \textbf{any mechanism must preserve $CPT$}.
All mechanisms listed above satisfy this.
\end{resultbox}

\ifdefined\INCLUDEMODE
\else
  \end{document}
\fi
